\documentclass[11pt]{article}

\usepackage[margin=1in]{geometry}
\usepackage{amsmath,amssymb,amsthm}
\usepackage{mathtools}
\usepackage{microtype}
\usepackage[hidelinks]{hyperref}

\hypersetup{
  pdftitle={Euclidean Support Modes of the Regular Dodecahedron},
  pdfauthor={Scott Allen Cave},
  pdfsubject={Euclidean convex geometry and dodecahedral support modes}
}

\setcounter{tocdepth}{1}
\setcounter{secnumdepth}{2}
\raggedbottom

\newtheorem{theorem}{Theorem}[section]
\newtheorem{proposition}[theorem]{Proposition}
\newtheorem{corollary}[theorem]{Corollary}
\newtheorem{lemma}[theorem]{Lemma}
\theoremstyle{definition}
\newtheorem{definition}[theorem]{Definition}
\theoremstyle{remark}
\newtheorem{remark}[theorem]{Remark}

\title{Euclidean Support Modes of the Regular Dodecahedron\\
\large Scale, Shape, and Symmetry from Six Opposite-Face Distances}
\author{Scott Allen Cave}
\date{July 11, 2026}

\begin{document}

\maketitle

\begin{abstract}
This paper develops a Euclidean scale-shape decomposition for a
six-parameter family based on the regular dodecahedron. Fix the twelve
regular dodecahedral face-normal directions and assign one positive support
number to each of the six opposite-face pairs. Their supporting slabs
intersect in an origin-centered, centrally symmetric convex polytope.

The resulting support register lies in \(\mathbb R^6\) and decomposes
orthogonally into a one-dimensional uniform subspace and a
five-dimensional zero-mean subspace. Common multiplication of all six
support numbers produces Euclidean homothety. While the defining planes
remain active facets, the uniform ray is the unique regular similarity
family, and every nonzero zero-mean residual produces a non-homothetic
support variation.

The decomposition is compatible with the rotational symmetry group of the
regular dodecahedron. Within the connected support chamber preserving
dodecahedral incidence, each nonregular state is described uniquely by mean
support, shape magnitude, and normalized shape direction. After removal of
absolute scale, the normalized directions lie on a four-sphere in the
five-dimensional zero-mean subspace.

The construction uses only Euclidean geometry, convex polytope theory,
finite-dimensional real linear algebra, and elementary symmetry actions.
It does not classify arbitrary dodecahedra and asserts no mechanical,
dynamical, or physical interpretation.
\end{abstract}

\tableofcontents

\section{Introduction}

The twelve faces of a regular dodecahedron occur in six opposite pairs.
Each pair determines an antipodal face-normal direction through the center
of the polyhedron. These six directions provide a natural coordinate system
for a restricted Euclidean realization family.

Fix the twelve face-normal directions and allow only the perpendicular
distances of the six opposite supporting plane pairs from the origin to
vary. Assigning one positive support number to each pair gives
\[
\mathbf h=(h_1,\ldots,h_6)\in\mathbb R_{>0}^6.
\]
The intersection of the corresponding symmetric slabs is an
origin-centered, centrally symmetric convex polytope.

The central question is:

\begin{quote}
How does the six-dimensional support register separate regular Euclidean
scale from relative shape variation?
\end{quote}

Define
\[
\bar h
=
\frac{1}{6}\sum_{i=1}^{6}h_i
\]
and
\[
\mathbf r
=
\mathbf h-\bar h\,\mathbf 1.
\]
Then
\[
\mathbf h
=
\bar h\,\mathbf 1+\mathbf r,
\qquad
\sum_{i=1}^{6}r_i=0.
\]

The algebraic decomposition is immediate. The main result of the paper is
its Euclidean interpretation.

Common multiplication of all support values scales the entire realization:
\[
P(c\mathbf h)=cP(\mathbf h).
\]
Conversely, while all defining planes remain active facets, a realization
homothetic to the regular reference dodecahedron must have six equal support
numbers. The uniform ray is therefore the unique regular scale direction,
and every nonzero residual gives a non-homothetic support variation.

The result is local to the connected chamber in support space that contains
the regular state and preserves dodecahedral incidence. Sufficiently small
support perturbations remain in this chamber by continuity and the strict
facet inequalities of the regular dodecahedron.

\subsection{Main results}

The paper establishes:

\begin{enumerate}
\item a complete six-coordinate description of the active fixed-normal,
origin-centered realization family;

\item the canonical orthogonal decomposition
\[
\mathbb R^6
=
\operatorname{span}\{\mathbf 1\}
\oplus
\left\{
\mathbf r\in\mathbb R^6:
\sum_i r_i=0
\right\};
\]

\item the identity
\[
P(c\mathbf h)=cP(\mathbf h),
\qquad
c>0;
\]

\item uniqueness of the regular scale direction while all defining planes
remain facet-defining;

\item local preservation of the dodecahedral combinatorial type; and

\item symmetry-compatible coordinates for mean scale, shape magnitude, and
normalized shape direction.
\end{enumerate}

\subsection{Scope}

All constructions take place in ordinary three-dimensional Euclidean space
and use standard convex polytope theory and finite-dimensional real linear
algebra.

The face-normal directions remain fixed, opposite supporting planes remain
equidistant from the origin, and central symmetry is preserved. Arbitrary
vertex displacement, independent face rotation, varying normal directions,
and non-centrally symmetric realizations are excluded.

No mechanics, elasticity, force, energy, dynamics, continuum model,
physical expansion, gravitation, or cosmological interpretation is
asserted.

\subsection{Organization}

Section~2 defines the support realization. Section~3 develops the canonical
split of the support register. Section~4 proves its Euclidean interpretation
as scale and shape. Section~5 records coordinate consequences, Section~6
gives examples, and Sections~7 and~8 discuss the scope and conclusions.

\section{Euclidean Setting and Support Coordinates}

This section defines the fixed-normal Euclidean family studied throughout
the paper.

\subsection{Ambient geometry}

Let
\[
\mathbb E^3
\]
denote three-dimensional Euclidean space with its standard inner product.

Let
\[
D\subset\mathbb E^3
\]
be a regular dodecahedron centered at the origin. Its twelve faces occur in
six opposite pairs. Throughout the paper, the twelve face-normal directions
of \(D\) remain fixed. Only the perpendicular distances of the corresponding
supporting planes from the origin are allowed to vary.

\subsection{Opposite-face normal pairs}

Choose one outward unit normal from each opposite-face pair:
\[
\mathcal N
=
\{\mathbf n_1,\ldots,\mathbf n_6\}.
\]

The complete set of outward normal directions is
\[
\{\pm\mathbf n_1,\ldots,\pm\mathbf n_6\}.
\]

The numbering is arbitrary. Every rotational symmetry of the regular
dodecahedron permutes these six antipodal normal pairs.

\subsection{Support realization}

Let
\[
\mathbf h=(h_1,\ldots,h_6)\in\mathbb R_{>0}^6.
\]

For each \(i\), the two supporting planes normal to
\(\pm\mathbf n_i\) are
\[
\mathbf n_i\cdot\mathbf x=\pm h_i.
\]

They bound the slab
\[
S_i(h_i)
=
\left\{
\mathbf x\in\mathbb E^3:
|\mathbf n_i\cdot\mathbf x|\le h_i
\right\}.
\]

\begin{definition}[Support realization]
For
\[
\mathbf h\in\mathbb R_{>0}^6,
\]
define
\[
P(\mathbf h)
=
\bigcap_{i=1}^{6}S_i(h_i)
=
\left\{
\mathbf x\in\mathbb E^3:
|\mathbf n_i\cdot\mathbf x|\le h_i
\text{ for all }i
\right\}.
\]
The vector \(\mathbf h\) is the \emph{support register}.
\end{definition}

The set \(P(\mathbf h)\) is convex as an intersection of closed
half-spaces. It is centrally symmetric because
\[
\mathbf x\in P(\mathbf h)
\quad\Longrightarrow\quad
-\mathbf x\in P(\mathbf h).
\]

The selected normal directions span \(\mathbb E^3\), so the intersection is
bounded for every positive support register. Thus \(P(\mathbf h)\) is a
centrally symmetric convex polytope.

\subsection{The regular state}

There exists \(h_0>0\) such that
\[
D=P(h_0\mathbf 1),
\qquad
\mathbf 1=(1,1,1,1,1,1).
\]

Hence the regular dodecahedron is represented by six equal support numbers.
The positive uniform ray is
\[
U_+
=
\{h\mathbf 1:h>0\}.
\]

Section~4 proves that this ray is precisely the regular Euclidean similarity
family in the fixed-normal realization.

\subsection{Support space and admissibility}

The support register lies in the six-dimensional vector space
\[
\mathbb R^6.
\]

The positive orthant
\[
\mathbb R_{>0}^6
\]
is the ambient support domain. Positivity keeps every opposite plane pair
at positive distance from the origin.

Positivity alone does not preserve dodecahedral incidence. Large support
variations may make a defining inequality redundant, contract an edge, or
alter the vertex structure. Section~4 therefore restricts the geometric
results to the connected support chamber containing \(h_0\mathbf 1\) and
preserving the dodecahedral combinatorial type.

\subsection{Completeness of the support register}

Within the active fixed-normal family, the support register uniquely
determines the realization.

\begin{proposition}[Uniqueness of the support register]
Let
\[
\mathbf h,\mathbf k\in\mathbb R_{>0}^6.
\]
Assume that all twelve defining planes support genuine facets of both
\(P(\mathbf h)\) and \(P(\mathbf k)\). If
\[
P(\mathbf h)=P(\mathbf k),
\]
then
\[
\mathbf h=\mathbf k.
\]
\end{proposition}

\begin{proof}
For each \(i\), the active facet with outward normal \(\mathbf n_i\) has
support value
\[
\max_{\mathbf x\in P(\mathbf h)}
\mathbf n_i\cdot\mathbf x
=
h_i.
\]

The same support value equals \(k_i\) for \(P(\mathbf k)\). Equal convex
polytopes have equal support values in every direction, so
\[
h_i=k_i
\]
for all \(i\).
\end{proof}

Thus the six coordinates are complete for the declared fixed-normal,
origin-centered, active-facet family.

\subsection{Symmetry action}

Let \(\Gamma\) be the rotational symmetry group of the regular
dodecahedron. Each \(g\in\Gamma\) induces a permutation
\[
\pi_g\in S_6
\]
of the opposite-face pairs.

Define
\[
\rho(g)(h_1,\ldots,h_6)
=
(h_{\pi_g^{-1}(1)},\ldots,h_{\pi_g^{-1}(6)}).
\]

Then
\[
\rho(g_1g_2)
=
\rho(g_1)\rho(g_2),
\]
so \(\rho\) is a linear permutation representation on \(\mathbb R^6\).

Each \(\rho(g)\) preserves the standard inner product, coordinate sums,
the positive orthant, and the uniform vector \(\mathbf 1\).

The realization is equivariant:
\[
P(\rho(g)\mathbf h)
=
gP(\mathbf h).
\]

Thus coordinate permutation in support space agrees with Euclidean rotation
of the realized polytope.

\subsection{Scope}

The family \(P(\mathbf h)\) varies only the six common distances of
opposite supporting plane pairs. It keeps fixed

\begin{itemize}
\item the twelve regular dodecahedral normal directions;
\item the common center at the origin;
\item equality of the two distances within each opposite pair; and
\item central symmetry.
\end{itemize}

It excludes varying normal directions, unequal opposite supports,
translations, independent face rotations, arbitrary vertex displacement,
and non-centrally symmetric realizations.

No mechanical or dynamical meaning is assigned to the support coordinates.
They record Euclidean plane positions only.

The next section derives the canonical decomposition of the support register
into its uniform and zero-mean components.

\section{The Canonical Split of the Support Register}

The support register
\[
\mathbf h=(h_1,\ldots,h_6)
\]
is a vector in \(\mathbb R^6\). This section separates it into a common
component shared by all six coordinates and a complementary zero-mean
component.

The decomposition is standard linear algebra. Its Euclidean meaning is
proved in Section~4.

\subsection{Uniform and zero-mean subspaces}

Let
\[
\mathbf 1=(1,1,1,1,1,1)
\]
and define the uniform subspace
\[
U=\operatorname{span}\{\mathbf 1\}.
\]

A register is uniform precisely when all six coordinates are equal.

Define also
\[
H=
\left\{
\mathbf r\in\mathbb R^6:
\sum_{i=1}^{6}r_i=0
\right\}.
\]

The subspace \(H\) is the kernel of the nonzero linear functional
\[
\mathbf x\longmapsto\sum_{i=1}^{6}x_i,
\]
and therefore
\[
\dim H=5.
\]

\subsection{Mean support and residual}

For any
\[
\mathbf h=(h_1,\ldots,h_6)\in\mathbb R^6,
\]
define its mean support by
\[
\bar h
=
\frac{1}{6}\sum_{i=1}^{6}h_i.
\]

The corresponding uniform component is
\[
\mathbf h_U
=
\bar h\,\mathbf 1.
\]

Define the residual by
\[
\mathbf r
=
\mathbf h-\bar h\,\mathbf 1.
\]

Its coordinate sum vanishes:
\[
\sum_{i=1}^{6}r_i
=
\sum_{i=1}^{6}h_i-6\bar h
=
0.
\]

Hence
\[
\mathbf r\in H.
\]

\subsection{Canonical decomposition}

\begin{theorem}[Canonical decomposition]
Every vector \(\mathbf h\in\mathbb R^6\) admits a unique decomposition
\[
\mathbf h
=
\bar h\,\mathbf 1+\mathbf r,
\]
where
\[
\bar h\,\mathbf 1\in U
\qquad\text{and}\qquad
\mathbf r\in H.
\]
\end{theorem}

\begin{proof}
The definitions above give existence.

For uniqueness, suppose
\[
\mathbf h
=
\mathbf u_1+\mathbf r_1
=
\mathbf u_2+\mathbf r_2,
\]
with
\[
\mathbf u_1,\mathbf u_2\in U
\qquad\text{and}\qquad
\mathbf r_1,\mathbf r_2\in H.
\]

Then
\[
\mathbf u_1-\mathbf u_2
=
\mathbf r_2-\mathbf r_1
\in U\cap H.
\]

If \(a\mathbf 1\in U\cap H\), then
\[
0=\sum_{i=1}^{6}a=6a,
\]
so \(a=0\). Thus
\[
U\cap H=\{\mathbf 0\},
\]
and both components are unique.
\end{proof}

Therefore
\[
\mathbb R^6=U\oplus H.
\]

\subsection{Orthogonality}

Equip \(\mathbb R^6\) with the standard inner product
\[
\langle\mathbf x,\mathbf y\rangle
=
\sum_{i=1}^{6}x_i y_i.
\]

For every \(\mathbf r\in H\),
\[
\langle\mathbf 1,\mathbf r\rangle
=
\sum_{i=1}^{6}r_i
=
0.
\]

Hence
\[
H=U^\perp,
\]
and the direct sum is orthogonal:
\[
\mathbb R^6
=
U\mathbin{\perp}H.
\]

In particular,
\[
\left\langle
\bar h\,\mathbf 1,
\mathbf r
\right\rangle
=
0.
\]

Thus
\[
\|\mathbf h\|^2
=
6\bar h^2+\|\mathbf r\|^2.
\]

\subsection{Projection operators}

Let \(J\) be the \(6\times6\) all-ones matrix. Define
\[
P_U=\frac{1}{6}J
\]
and
\[
P_H=I-\frac{1}{6}J.
\]

Then
\[
P_U\mathbf h
=
\bar h\,\mathbf 1
\]
and
\[
P_H\mathbf h
=
\mathbf r.
\]

Since
\[
J^2=6J,
\]
the operators satisfy
\[
P_U^2=P_U,
\qquad
P_H^2=P_H,
\]
\[
P_UP_H=P_HP_U=0,
\]
and
\[
P_U+P_H=I.
\]

They are therefore the complementary orthogonal projections onto \(U\)
and \(H\).

\subsection{Symmetry compatibility}

Every rotational symmetry of the regular dodecahedron acts on support space
by a permutation matrix \(\rho(g)\).

Coordinate permutations fix \(\mathbf 1\) and preserve coordinate sums.
Therefore
\[
\rho(g)U=U
\qquad\text{and}\qquad
\rho(g)H=H.
\]

The projection operators commute with the symmetry action:
\[
\rho(g)P_U=P_U\rho(g),
\qquad
\rho(g)P_H=P_H\rho(g).
\]

Consequently,
\[
P_U(\rho(g)\mathbf h)
=
\rho(g)(P_U\mathbf h)
\]
and
\[
P_H(\rho(g)\mathbf h)
=
\rho(g)(P_H\mathbf h).
\]

The split is therefore independent of the arbitrary labeling of the six
opposite-face pairs.

\subsection{Positive support registers}

The algebraic decomposition applies to every vector in \(\mathbb R^6\).
The geometric support realization requires
\[
\mathbf h\in\mathbb R_{>0}^6.
\]

If
\[
\mathbf h
=
\bar h\,\mathbf 1+\mathbf r
\]
is positive, then
\[
\bar h+r_i>0
\]
for every \(i\), and necessarily
\[
\bar h>0.
\]

The residual may contain positive and negative entries. These are deviations
from a positive mean support, not negative support distances.

\subsection{Interpretation boundary}

The canonical split establishes one uniform coordinate and five independent
zero-mean coordinates.

It does not by itself prove that the uniform component is Euclidean scale
or that the residual is geometric shape variation.

Section~4 proves those conclusions from the support realization.

\section{Euclidean Scale and Shape}

Section~3 established the orthogonal decomposition
\[
\mathbf h
=
\bar h\,\mathbf 1+\mathbf r,
\qquad
\sum_{i=1}^{6}r_i=0.
\]

This section proves its geometric meaning for the fixed-normal support
realization
\[
P(\mathbf h)
=
\left\{
\mathbf x\in\mathbb E^3:
|\mathbf n_i\cdot\mathbf x|\le h_i
\text{ for all }i
\right\}.
\]

The uniform direction is exactly the regular Euclidean scale direction.
Within the dodecahedral support chamber, a nonzero residual gives a
non-homothetic shape variation.

\subsection{Support scaling}

\begin{theorem}[Support scaling]
For every
\[
\mathbf h\in\mathbb R_{>0}^6
\]
and every \(c>0\),
\[
P(c\mathbf h)=cP(\mathbf h).
\]
\end{theorem}

\begin{proof}
Let \(\mathbf x\in P(c\mathbf h)\) and set
\[
\mathbf y=\frac{1}{c}\mathbf x.
\]
Then
\[
|\mathbf n_i\cdot\mathbf y|
=
\frac{1}{c}|\mathbf n_i\cdot\mathbf x|
\le h_i
\]
for every \(i\). Hence
\[
\mathbf y\in P(\mathbf h),
\]
so
\[
\mathbf x\in cP(\mathbf h).
\]

Conversely, if
\[
\mathbf x=c\mathbf y
\]
with \(\mathbf y\in P(\mathbf h)\), then
\[
|\mathbf n_i\cdot\mathbf x|
=
c|\mathbf n_i\cdot\mathbf y|
\le ch_i.
\]
Thus
\[
\mathbf x\in P(c\mathbf h).
\]
\end{proof}

This result holds throughout the positive support domain and does not
require preservation of dodecahedral incidence.

\subsection{The regular similarity family}

Let
\[
D=P(h_0\mathbf 1)
\]
be the reference regular dodecahedron.

For every \(h>0\),
\[
P(h\mathbf 1)
=
\frac{h}{h_0}D.
\]

Therefore the positive uniform ray
\[
U_+
=
\{h\mathbf 1:h>0\}
\]
realizes the regular origin-centered similarity family with the fixed
normal directions.

\subsection{Uniqueness of the regular scale direction}

For a compact convex set \(K\subset\mathbb E^3\), define its support value
in the unit direction \(\mathbf u\) by
\[
s_K(\mathbf u)
=
\max_{\mathbf x\in K}
\mathbf u\cdot\mathbf x.
\]

If the plane
\[
\mathbf n_i\cdot\mathbf x=h_i
\]
supports a genuine facet of \(P(\mathbf h)\), then
\[
s_{P(\mathbf h)}(\mathbf n_i)=h_i.
\]

\begin{theorem}[Uniqueness of the regular scale direction]
Let
\[
\mathbf h\in\mathbb R_{>0}^6,
\]
and assume that all twelve defining planes of \(P(\mathbf h)\) support
genuine facets. If
\[
P(\mathbf h)=cD
\]
for some \(c>0\), then
\[
\mathbf h=ch_0\mathbf 1.
\]
\end{theorem}

\begin{proof}
For every \(i\),
\[
h_i
=
s_{P(\mathbf h)}(\mathbf n_i).
\]

Since support values scale under homothety,
\[
s_{P(\mathbf h)}(\mathbf n_i)
=
c\,s_D(\mathbf n_i)
=
ch_0.
\]

Thus
\[
h_i=ch_0
\]
for all \(i\), and hence
\[
\mathbf h=ch_0\mathbf 1.
\]
\end{proof}

The active-facet assumption excludes redundant inequalities whose written
support values are not recoverable from the realized polytope.

\subsection{Residual variation}

Write
\[
\mathbf h
=
\bar h\,\mathbf 1+\mathbf r,
\qquad
\mathbf r\in H.
\]

If
\[
\mathbf r=\mathbf 0,
\]
then
\[
\mathbf h=\bar h\,\mathbf 1,
\]
and \(P(\mathbf h)\) is regular.

If
\[
\mathbf r\ne\mathbf 0,
\]
then the six support numbers are not all equal. Under the active-facet
condition, the uniqueness theorem shows that \(P(\mathbf h)\) cannot be a
homothetic copy of the regular dodecahedron.

\begin{corollary}[Residual criterion]
Assume that all twelve defining planes remain facet-defining. Then
\[
\mathbf r=\mathbf 0
\]
if and only if \(P(\mathbf h)\) lies in the regular similarity family.
\end{corollary}

\subsection{Local preservation of dodecahedral incidence}

Positive support numbers do not alone guarantee the regular dodecahedral
combinatorial type. The required stability holds locally.

\begin{theorem}[Local stability of dodecahedral incidence]
There exists an open neighborhood
\[
\mathcal U\subset\mathbb R_{>0}^6
\]
of \(h_0\mathbf 1\) such that every \(P(\mathbf h)\) with
\[
\mathbf h\in\mathcal U
\]
has the same face-vertex incidence structure as the regular dodecahedron.
\end{theorem}

\begin{proof}
Write the twelve signed unit normals as
\[
\mathbf m_1,\ldots,\mathbf m_{12},
\]
and let \(p(a)\in\{1,\ldots,6\}\) identify the opposite-face pair
containing \(\mathbf m_a\).

Then
\[
P(\mathbf h)
=
\left\{
\mathbf x:
\mathbf m_a\cdot\mathbf x
\le
h_{p(a)}
\text{ for }a=1,\ldots,12
\right\}.
\]

The regular dodecahedron is simple, so each vertex is determined by an
invertible triple
\[
T=\{a,b,c\}
\]
of incident facet planes. Let
\[
N_T
=
\begin{pmatrix}
\mathbf m_a^{\mathsf T}\\
\mathbf m_b^{\mathsf T}\\
\mathbf m_c^{\mathsf T}
\end{pmatrix}.
\]

The corresponding intersection point is
\[
\mathbf v_T(\mathbf h)
=
N_T^{-1}
\begin{pmatrix}
h_{p(a)}\\
h_{p(b)}\\
h_{p(c)}
\end{pmatrix},
\]
which depends continuously on \(\mathbf h\).

At the regular state, every nonincident facet inequality is strict at each
regular vertex. Because only finitely many such inequalities occur, their
minimum slack is positive. Hence all regular vertex triples remain feasible
throughout a sufficiently small neighborhood.

Conversely, every invertible triple that is not a regular vertex violates
at least one defining inequality strictly at the regular state. Finiteness
and continuity allow the neighborhood to be chosen so that no such triple
becomes feasible.

After also preserving positivity, exactly the original vertex triples
remain feasible. Therefore the vertex-facet incidence structure, and hence
the full dodecahedral combinatorial type, is unchanged.
\end{proof}

\subsection{The dodecahedral support chamber}

\begin{definition}[Dodecahedral support chamber]
Let \(\mathcal C_D\) be the connected component containing
\(h_0\mathbf 1\) of the set of positive support registers for which

\begin{enumerate}
\item all twelve defining planes support genuine facets;
\item the polytope is simple; and
\item its face-vertex incidence structure is dodecahedral.
\end{enumerate}
\]
\end{definition}

The local stability theorem guarantees that \(\mathcal C_D\) contains an
open neighborhood of the regular state.

No claim is made that
\[
\mathcal C_D=\mathbb R_{>0}^6.
\]

Its boundary may contain facet loss, edge contraction, vertex collision, or
other changes of combinatorial type.

\subsection{Euclidean scale-shape theorem}

\begin{theorem}[Euclidean scale-shape decomposition]
Let
\[
\mathbf h\in\mathcal C_D
\]
and write
\[
\mathbf h
=
\bar h\,\mathbf 1+\mathbf r,
\qquad
\sum_{i=1}^{6}r_i=0.
\]

Then:

\begin{enumerate}
\item \(\bar h\,\mathbf 1\) is the unique regular support register with mean
support \(\bar h\);

\item for every \(c>0\),
\[
P(c\mathbf h)=cP(\mathbf h);
\]

\item \(P(\mathbf h)\) is regular if and only if
\[
\mathbf r=\mathbf 0;
\]

\item if
\[
\mathbf r\ne\mathbf 0,
\]
then \(P(\mathbf h)\) is not an origin-centered homothetic copy of \(D\).
\end{enumerate}
\end{theorem}

\begin{proof}
The first statement follows from the regular uniform family and uniqueness
of the regular scale direction. The second is the support-scaling theorem.

If \(\mathbf r=\mathbf 0\), then the support register is uniform and the
realization is regular. Conversely, a regular realization in
\(\mathcal C_D\) must have equal support values, so its residual vanishes.
The final statement follows immediately.
\end{proof}

Thus
\[
\text{uniform component}
\longleftrightarrow
\text{regular Euclidean scale},
\]
while
\[
\text{nonzero zero-mean residual}
\longleftrightarrow
\text{non-homothetic support shape variation}.
\]

\subsection{Boundary}

The theorem applies only to the fixed-normal, origin-centered, centrally
symmetric realization family and only inside \(\mathcal C_D\).

It does not classify arbitrary Euclidean dodecahedra or arbitrary
deformations of the regular solid. It excludes varying normal directions,
unequal opposite supports, translations, independent face rotations, and
arbitrary vertex motion.

The terms \emph{scale} and \emph{shape} are used only in their ordinary
Euclidean sense. No mechanics, elasticity, force, energy, dynamics, or
physical expansion is implied.

\section{Coordinate Consequences}

The Euclidean scale-shape theorem identifies the uniform component of the
support register with regular scale and the zero-mean residual with
non-homothetic support variation. This section records the resulting
coordinates.

Let
\[
\mathbf h
=
\bar h\,\mathbf 1+\mathbf r
\in\mathcal C_D,
\qquad
\sum_{i=1}^{6}r_i=0.
\]

\subsection{Mean support and residual}

The mean support is
\[
\bar h
=
\frac{1}{6}\sum_{i=1}^{6}h_i.
\]

Relative to the reference regular dodecahedron
\[
D=P(h_0\mathbf 1),
\]
the associated regular scale ratio is
\[
\lambda(\mathbf h)
=
\frac{\bar h}{h_0}.
\]

Hence
\[
P(\bar h\,\mathbf 1)
=
\lambda(\mathbf h)D.
\]

The shape residual is
\[
\mathbf r
=
\mathbf h-\bar h\,\mathbf 1,
\]
with coordinates
\[
r_i=h_i-\bar h.
\]

Adding a common amount to every support coordinate leaves the residual
unchanged:
\[
(\mathbf h+a\mathbf 1)
-
\overline{(\mathbf h+a\mathbf 1)}\,\mathbf 1
=
\mathbf r.
\]

Under positive scaling,
\[
\mathbf h\longmapsto c\mathbf h,
\]
one has
\[
\bar h\longmapsto c\bar h
\qquad\text{and}\qquad
\mathbf r\longmapsto c\mathbf r.
\]

\subsection{Shape magnitude}

Define
\[
\sigma(\mathbf h)
=
\|\mathbf r\|
=
\left(
\sum_{i=1}^{6}(h_i-\bar h)^2
\right)^{1/2}.
\]

This is the Euclidean distance from \(\mathbf h\) to the uniform subspace:
\[
\sigma(\mathbf h)
=
\operatorname{dist}(\mathbf h,U).
\]

\begin{corollary}[Regularity criterion]
For every \(\mathbf h\in\mathcal C_D\),
\[
\sigma(\mathbf h)=0
\]
if and only if \(P(\mathbf h)\) is regular.
\end{corollary}

The quantity \(\sigma\) measures support-space departure from the regular
ray. It is not asserted to be a mechanical strain, energy, or physical
deformation cost.

\subsection{Shape direction}

When
\[
\sigma(\mathbf h)>0,
\]
define the normalized residual
\[
\widehat{\mathbf r}
=
\frac{\mathbf r}{\|\mathbf r\|}.
\]

Then
\[
\widehat{\mathbf r}\in H,
\qquad
\|\widehat{\mathbf r}\|=1.
\]

Since \(\dim H=5\), normalized shape directions lie on
\[
S(H)
=
\left\{
\mathbf v\in H:
\|\mathbf v\|=1
\right\}
\cong S^4.
\]

Every nonregular support register therefore has the unique form
\[
\mathbf h
=
\bar h\,\mathbf 1
+
\sigma\widehat{\mathbf r}.
\]

Thus a nonregular state is specified by

\begin{enumerate}
\item mean support \(\bar h\);
\item shape magnitude \(\sigma\); and
\item normalized shape direction \(\widehat{\mathbf r}\in S^4\).
\end{enumerate}

At a regular state, \(\sigma=0\), and no shape direction is defined.

\subsection{Scale-free shape}

Because
\[
\sigma(c\mathbf h)=c\sigma(\mathbf h),
\]
the magnitude \(\sigma\) depends on absolute scale.

Define the relative shape magnitude
\[
\eta(\mathbf h)
=
\frac{\sigma(\mathbf h)}{\bar h}.
\]

Since \(\bar h>0\) on the positive support domain, \(\eta\) is well defined.
It is invariant under positive scaling:
\[
\eta(c\mathbf h)
=
\eta(\mathbf h).
\]

Define also the normalized support register
\[
\mathbf q
=
\frac{\mathbf h}{\bar h}.
\]

Then
\[
\frac{1}{6}\sum_{i=1}^{6}q_i=1
\]
and
\[
\mathbf q
=
\mathbf 1+\eta\widehat{\mathbf r}.
\]

Equivalently,
\[
\mathbf h
=
\bar h
\left(
\mathbf 1+\eta\widehat{\mathbf r}
\right).
\]

The factor \(\bar h\) records absolute scale. The pair
\[
(\eta,\widehat{\mathbf r})
\]
records scale-free support shape.

\subsection{Orthogonal identity}

The canonical decomposition is orthogonal, so
\[
\|\mathbf h\|^2
=
\|\bar h\,\mathbf 1\|^2+\|\mathbf r\|^2.
\]

Since
\[
\|\mathbf 1\|^2=6,
\]
this becomes
\[
\|\mathbf h\|^2
=
6\bar h^2+\sigma^2.
\]

After division by \(\bar h^2\),
\[
\frac{\|\mathbf h\|^2}{\bar h^2}
=
6+\eta^2.
\]

These are identities in the Euclidean metric on support space. They are not
formulas for volume, surface area, energy, or any other physical quantity.

\subsection{Symmetry behavior}

For every rotational symmetry \(g\),
\[
\overline{\rho(g)\mathbf h}
=
\bar h.
\]

Because \(\rho(g)\) is orthogonal,
\[
\sigma(\rho(g)\mathbf h)
=
\sigma(\mathbf h)
\]
and
\[
\eta(\rho(g)\mathbf h)
=
\eta(\mathbf h).
\]

The normalized direction transforms equivariantly:
\[
\widehat{\mathbf r}(\rho(g)\mathbf h)
=
\rho(g)\widehat{\mathbf r}(\mathbf h).
\]

Thus \(\bar h\), \(\sigma\), and \(\eta\) are scalar symmetry invariants,
whereas \(\widehat{\mathbf r}\) retains the oriented distribution of support
variation among the six face pairs.

\subsection{Shape space and basis choice}

The five-dimensional subspace
\[
H
=
\left\{
\mathbf r\in\mathbb R^6:
\sum_i r_i=0
\right\}
\]
is canonical. An ordered basis within it is not.

For example,
\[
\begin{aligned}
\mathbf b_1&=(1,-1,0,0,0,0),\\
\mathbf b_2&=(0,1,-1,0,0,0),\\
\mathbf b_3&=(0,0,1,-1,0,0),\\
\mathbf b_4&=(0,0,0,1,-1,0),\\
\mathbf b_5&=(0,0,0,0,1,-1)
\end{aligned}
\]
span \(H\), but depend on the chosen ordering of the opposite-face pairs.

A preferred basis would require an additional criterion. No such choice is
made here.

\subsection{Summary}

For every support register in \(\mathcal C_D\),
\[
\mathbf h
=
\bar h\,\mathbf 1+\mathbf r.
\]

For nonregular states,
\[
\mathbf h
=
\bar h\,\mathbf 1
+
\sigma\widehat{\mathbf r}
=
\bar h
\left(
\mathbf 1+\eta\widehat{\mathbf r}
\right).
\]

The resulting coordinates separate

\begin{itemize}
\item absolute scale through \(\bar h\);
\item absolute shape departure through \(\sigma\);
\item scale-free shape departure through \(\eta\); and
\item shape direction through \(\widehat{\mathbf r}\in S^4\).
\end{itemize}

Section~6 applies these coordinates to representative support registers.

\section{Representative Support Modes}

This section applies the scale-shape coordinates to several explicit support
registers. The examples are representatives only. They do not define a
canonical basis of the five-dimensional shape space.

Let
\[
\mathbf h_0=h_0\mathbf 1,
\qquad
h_0>0,
\]
be the reference regular support register. Every perturbed register below is
assumed to remain in the dodecahedral support chamber \(\mathcal C_D\).

\subsection{Regular state and uniform scale}

For the reference state
\[
\mathbf h=h_0\mathbf 1,
\]
one has
\[
\bar h=h_0,
\qquad
\mathbf r=\mathbf 0,
\qquad
\sigma=0,
\qquad
\eta=0.
\]

The realization is the regular dodecahedron:
\[
P(\mathbf h)=D.
\]

More generally, let
\[
\mathbf h=ch_0\mathbf 1,
\qquad
c>0.
\]

Then
\[
\bar h=ch_0,
\qquad
\mathbf r=\mathbf 0,
\qquad
\sigma=\eta=0,
\]
and
\[
P(\mathbf h)=cD.
\]

Thus the positive uniform ray changes only regular Euclidean scale.

\subsection{Balanced exchange between two face pairs}

Let
\[
\mathbf v_1=(1,-1,0,0,0,0)
\]
and consider
\[
\mathbf h
=
h_0\mathbf 1+a\mathbf v_1.
\]

Since \(\mathbf v_1\) has zero coordinate sum,
\[
\bar h=h_0.
\]

The residual is
\[
\mathbf r
=
a(1,-1,0,0,0,0),
\]
with magnitude
\[
\sigma
=
\sqrt{2}\,|a|.
\]

For \(a\ne0\),
\[
\widehat{\mathbf r}
=
\operatorname{sgn}(a)
\frac{1}{\sqrt{2}}
(1,-1,0,0,0,0),
\]
and
\[
\eta
=
\frac{\sqrt{2}\,|a|}{h_0}.
\]

One opposite-face pair moves outward relative to the mean, one moves inward,
and four remain unchanged. Since the residual is nonzero, the realization is
not homothetic to \(D\).

\subsection{One dominant face pair}

Let
\[
\mathbf v_2
=
(5,-1,-1,-1,-1,-1)
\]
and consider
\[
\mathbf h
=
h_0\mathbf 1+a\mathbf v_2.
\]

Again,
\[
\bar h=h_0.
\]

The residual satisfies
\[
\mathbf r
=
a(5,-1,-1,-1,-1,-1),
\]
and
\[
\sigma
=
\sqrt{30}\,|a|.
\]

For \(a\ne0\),
\[
\widehat{\mathbf r}
=
\operatorname{sgn}(a)
\frac{1}{\sqrt{30}}
(5,-1,-1,-1,-1,-1),
\]
with
\[
\eta
=
\frac{\sqrt{30}\,|a|}{h_0}.
\]

This register distinguishes one opposite-face pair from the remaining five.
A different dominant pair gives a symmetry-related example.

\subsection{A four-pair balanced mode}

Let
\[
\mathbf v_3
=
(1,1,-1,-1,0,0)
\]
and consider
\[
\mathbf h
=
h_0\mathbf 1+a\mathbf v_3.
\]

Then
\[
\bar h=h_0,
\qquad
\mathbf r
=
a(1,1,-1,-1,0,0),
\]
and
\[
\sigma=2|a|.
\]

For \(a\ne0\),
\[
\widehat{\mathbf r}
=
\operatorname{sgn}(a)
\frac{1}{2}
(1,1,-1,-1,0,0),
\]
while
\[
\eta
=
\frac{2|a|}{h_0}.
\]

Two face pairs move outward relative to the mean, two move inward, and two
remain at the mean support.

\subsection{Mixed scale and shape}

Consider
\[
\mathbf h
=
(h_0+\delta)\mathbf 1
+
a(1,-1,0,0,0,0),
\]
with
\[
h_0+\delta>0.
\]

Its mean support is
\[
\bar h=h_0+\delta,
\]
while its residual is
\[
\mathbf r
=
a(1,-1,0,0,0,0).
\]

Therefore
\[
\sigma
=
\sqrt{2}\,|a|
\]
and
\[
\eta
=
\frac{\sqrt{2}\,|a|}{h_0+\delta}.
\]

The same absolute residual has smaller relative magnitude when the mean
support is larger. This separates the roles of \(\sigma\) and \(\eta\).

\subsection{Normalized support registers}

For the balanced two-pair example,
\[
\mathbf h
=
h_0\mathbf 1+a(1,-1,0,0,0,0),
\]
the normalized register is
\[
\mathbf q
=
\frac{\mathbf h}{h_0}
=
\left(
1+\frac{a}{h_0},
1-\frac{a}{h_0},
1,1,1,1
\right).
\]

Its mean coordinate is one.

In general,
\[
\mathbf q
=
\frac{\mathbf h}{\bar h}
=
\mathbf 1+\eta\widehat{\mathbf r}.
\]

Thus normalized registers lie in the affine hyperplane
\[
\left\{
\mathbf q\in\mathbb R^6:
\frac{1}{6}\sum_i q_i=1
\right\}.
\]

The regular shape is represented by \(\mathbf 1\). Positive scalar multiples
of the same support register have the same normalized shape coordinate.

\subsection{Symmetry-related registers}

If \(g\) is a rotational symmetry and
\[
\mathbf h'=\rho(g)\mathbf h,
\]
then
\[
P(\mathbf h')=gP(\mathbf h).
\]

The two realizations are congruent, with permuted face-pair labels. They
satisfy
\[
\bar h(\mathbf h')=\bar h(\mathbf h),
\]
\[
\sigma(\mathbf h')=\sigma(\mathbf h),
\]
and
\[
\eta(\mathbf h')=\eta(\mathbf h).
\]

Their normalized residuals are related by
\[
\widehat{\mathbf r}(\mathbf h')
=
\rho(g)\widehat{\mathbf r}(\mathbf h).
\]

Examples in the same symmetry orbit therefore describe the same unlabeled
Euclidean shape.

\subsection{Admissibility}

The algebraic formulas are valid whenever the support coordinates are
positive. Their interpretation as dodecahedral modes requires
\[
\mathbf h\in\mathcal C_D.
\]

The local stability theorem guarantees this for sufficiently small
perturbations of \(h_0\mathbf 1\). The paper does not compute a global
maximum for \(a\) or \(\delta\).

At the boundary of \(\mathcal C_D\), a facet may become redundant, an edge
may contract, vertices may merge, or the incidence structure may change.
Such states lie outside the admitted dodecahedral mode family.

\subsection{Summary}

The examples display three basic cases:

\begin{enumerate}
\item uniform registers change only scale;
\item zero-mean registers change shape at fixed mean support; and
\item mixed registers change both scale and shape.
\end{enumerate}

For each register,
\[
\bar h
\]
records mean scale,
\[
\sigma
\]
records absolute shape departure,
\[
\eta
\]
records scale-free shape departure, and
\[
\widehat{\mathbf r}\in S^4
\]
records shape direction.

No mechanical or dynamical interpretation is introduced.

\section{Discussion}

The paper asked how far classical Euclidean and convex geometry could
separate scale from shape in a finite dodecahedral realization family.

Within the fixed-normal, origin-centered family, the answer is exact. The
six support numbers determine the active realization, their uniform
component gives the regular scale direction, and their zero-mean component
gives non-homothetic support variation.

\subsection{Canonical structure}

The six opposite-face pairs determine the support coordinates. Their common
direction
\[
U=\operatorname{span}\{\mathbf 1\}
\]
and zero-mean complement
\[
H=
\left\{
\mathbf r\in\mathbb R^6:
\sum_i r_i=0
\right\}
\]
are independent of the labeling of those pairs.

The projection operators
\[
P_U=\frac{1}{6}J
\qquad\text{and}\qquad
P_H=I-\frac{1}{6}J
\]
are therefore canonical.

The scalar quantities
\[
\bar h,
\qquad
\sigma,
\qquad
\eta
\]
are invariant under the rotational symmetry group of the regular
dodecahedron. The normalized residual
\[
\widehat{\mathbf r}
\]
is equivariant: a rotation permutes its coordinates in the same way that it
permutes the opposite-face pairs.

\subsection{Basis choice}

The shape space \(H\) is canonical, but an ordered basis within it is not.

Vectors such as
\[
(1,-1,0,0,0,0)
\]
and
\[
(5,-1,-1,-1,-1,-1)
\]
are useful representatives, but their coordinate forms depend on the
chosen labeling.

A preferred basis would require an additional criterion, such as a selected
subgroup, a symmetry-compatible operator, or an extremal geometric
quantity. No such criterion is introduced here.

\subsection{Role of the fixed-normal assumption}

The fixed-normal assumption makes the realization family finite and
explicit. It also makes the scale direction exact, because common
multiplication of the six support values produces Euclidean homothety.

The same assumption excludes many ordinary Euclidean deformations. A more
general dodecahedron may rotate face normals, break central symmetry, alter
opposition, or move vertices independently.

The present results therefore describe one natural six-parameter family,
not all Euclidean dodecahedra.

\subsection{Support chamber}

The chamber
\[
\mathcal C_D
\]
is the connected region in which all defining planes remain active and the
dodecahedral incidence structure is preserved.

Its local existence follows from finite strict slack and continuity. The
paper does not determine its global boundary.

A full chamber computation would identify the support registers at which a
facet becomes redundant, an edge contracts, a vertex triple loses
feasibility, a new vertex triple becomes feasible, or simplicity fails.

That is a finite problem in classical polyhedral geometry, but it is not
needed for the local scale-shape theorem.

\subsection{Meaning of scale and shape}

Scale means positive homothety about the origin:
\[
P(\mathbf h)
\longmapsto
cP(\mathbf h).
\]

Shape variation means departure from the regular uniform ray within the
fixed-normal chamber.

The quantity
\[
\sigma=\|\mathbf r\|
\]
is the Euclidean distance from the support register to the uniform
subspace. The ratio
\[
\eta=\frac{\sigma}{\bar h}
\]
removes common scale.

These quantities provide useful coordinates, but they are not claimed to be
the only possible measures of geometric shape.

\subsection{Classical sufficiency}

No new geometric primitive is required.

The construction uses only

\begin{itemize}
\item Euclidean planes, distances, and similarities;
\item support values of convex polytopes;
\item finite intersections of half-spaces;
\item orthogonal projection in \(\mathbb R^6\);
\item continuity of finite plane intersections; and
\item permutation actions induced by Euclidean symmetry.
\end{itemize}

These tools supply

\begin{itemize}
\item a complete fixed-normal coordinate register;
\item a unique regular scale direction;
\item a five-dimensional shape space;
\item a local chamber preserving dodecahedral incidence; and
\item symmetry-compatible scale-free shape coordinates.
\end{itemize}

The result therefore demonstrates classical mathematical sufficiency for
the declared problem.

\subsection{Algebra and realization}

The decomposition
\[
\mathbb R^6=U\oplus H
\]
is elementary. Its geometric content comes from the realization map
\[
\mathbf h\longmapsto P(\mathbf h).
\]

Without that map, the split separates only a mean coordinate from zero-mean
variation.

With it,
\[
U_+
\longleftrightarrow
\text{regular homothetic dodecahedra},
\]
while nonzero residuals in the active chamber correspond to
non-homothetic support variation.

The contribution is therefore the alignment of a familiar linear
decomposition with an explicit Euclidean realization.

\subsection{Limitations}

The paper does not provide

\begin{itemize}
\item a classification of all Euclidean dodecahedra;
\item a global description of \(\mathcal C_D\);
\item a canonical basis of \(H\);
\item a mechanical deformation theory;
\item a dynamical or constitutive law;
\item a force or energy model;
\item a continuum limit; or
\item a physical expansion law.
\end{itemize}

Every conclusion ends at the fixed-normal support family.

\subsection{Further classical questions}

Several extensions remain within ordinary geometry.

One may compute the full support chamber by enumerating the feasibility
conditions of candidate vertex triples.

One may study volume, surface area, edge lengths, and dihedral angles as
functions of the support register.

One may analyze the five-dimensional shape representation under selected
subgroups or additional symmetry-compatible operators.

One may also enlarge the realization family by allowing the face-normal
directions to vary, though that would require additional coordinates and a
new admissibility analysis.

\subsection{Summary}

The regular dodecahedron supplies six opposite-face support coordinates.
Convex geometry turns them into a complete fixed-normal realization, and
linear algebra separates their common mean from their relative variation.

Within \(\mathcal C_D\), the common direction is exactly regular Euclidean
scale, and the complementary directions are non-homothetic support shape
variation.

The result is local, bounded, and classical. Its value lies in showing how
much structure ordinary Euclidean geometry already provides.

\section{Conclusion}

This paper developed a Euclidean support-coordinate description of scale
and shape for a fixed-normal family based on the regular dodecahedron.

The six pairs of opposite face-normal directions define the support
register
\[
\mathbf h=(h_1,\ldots,h_6)\in\mathbb R_{>0}^6.
\]

Its realization is
\[
P(\mathbf h)
=
\left\{
\mathbf x\in\mathbb E^3:
|\mathbf n_i\cdot\mathbf x|\le h_i
\text{ for }i=1,\ldots,6
\right\}.
\]

Within the active fixed-normal family, these six coordinates uniquely
determine the polytope.

The support space admits the canonical orthogonal decomposition
\[
\mathbb R^6=U\oplus H,
\]
where
\[
U=\operatorname{span}\{\mathbf 1\}
\]
and
\[
H=
\left\{
\mathbf r\in\mathbb R^6:
\sum_{i=1}^{6}r_i=0
\right\}.
\]

Accordingly,
\[
\mathbf h
=
\bar h\,\mathbf 1+\mathbf r,
\qquad
\bar h
=
\frac{1}{6}\sum_{i=1}^{6}h_i.
\]

The geometric meaning of this split is exact.

For every \(c>0\),
\[
P(c\mathbf h)=cP(\mathbf h).
\]

Thus the positive uniform ray gives the regular similarity family. While all
defining planes remain active facets, a realization homothetic to the
regular reference dodecahedron must have equal support values. The uniform
ray is therefore the unique regular scale direction.

Within the dodecahedral support chamber,
\[
\mathbf r=\mathbf 0
\]
if and only if the realization is regular. Every nonzero residual produces
a non-homothetic support variation.

The decomposition is preserved by the rotational symmetry group. Its
orthogonal projections are
\[
P_U=\frac{1}{6}J
\]
and
\[
P_H=I-\frac{1}{6}J.
\]

Every nonregular support state has the unique form
\[
\mathbf h
=
\bar h\,\mathbf 1
+
\sigma\widehat{\mathbf r},
\]
where
\[
\sigma=\|\mathbf r\|>0
\]
and
\[
\widehat{\mathbf r}\in S^4\subset H.
\]

Equivalently,
\[
\mathbf h
=
\bar h
\left(
\mathbf 1+\eta\widehat{\mathbf r}
\right),
\qquad
\eta=\frac{\sigma}{\bar h}.
\]

The admitted family therefore carries one scale coordinate and five shape
degrees of freedom, organized by magnitude and direction.

The result is bounded to origin-centered, centrally symmetric realizations
with the regular dodecahedral normal directions fixed and the dodecahedral
incidence structure preserved. It does not classify arbitrary Euclidean
dodecahedra or arbitrary deformations of the regular solid.

No mechanics, elasticity, constitutive law, force, energy, dynamics,
continuum model, physical expansion, gravitation, or cosmological
interpretation has been introduced.

The main conclusion is classical:

\begin{quote}
For the fixed-normal dodecahedral support family, Euclidean convex geometry
and finite-dimensional real linear algebra are sufficient to separate
regular scale from relative shape.
\end{quote}

The construction does not require mathematics to be extended. It requires
the available mathematics to be used carefully and carried to its natural
boundary.


\end{document}
